\section{From Hessian bounds to triangle comparison}
\label{sec:weighted-comparison}

The preceding section bounds the distributional Hessian of squared
distance using the sectional-curvature assumption. We now turn that
bound into a differential inequality along every geodesic and then
into triangle comparison. The only new analytic input is the general
weak-Hessian theorem proved in the companion article
\cite[Theorem~1.1]{WeakHessianCompanion}. We record its exact interface
before constructing the distance functions to which it applies.

\Needspace{20\baselineskip}
\subsection{The general weak-Hessian input}

\begin{theorem}[Weak Hessian bounds along every geodesic]
\label{weight:every-geodesic}
Let \((M,d,m)\) be a full-support \(\RCD(K,N)\) space, where
\(K\in\mathbb R\) and \(1<N<\infty\). Let \(F\) be bounded and
globally Lipschitz, and let \(G\) be bounded and continuous.
Suppose that for every compactly supported \(g\in\Test(M)\) and
every nonnegative \(h\in\Lip_c(M)\),
\begin{equation}
\label{weight:hessian-assumption}
 H_F(\nabla g,\nabla g)(h)
 \le \int hG|\nabla g|^2\dd m,
\end{equation}
where the weak Hessian is defined by \eqref{eq:weak-hessian}.
Then every constant-speed minimizing geodesic
\(\sigma:[0,1]\to M\), of length \(\ell\), satisfies
\begin{equation}
\label{weight:curve-hessian}
 (F\circ\sigma)''\le\ell^2G\circ\sigma
 \quad\text{in distributions on }(0,1).
\end{equation}
\end{theorem}

This is \cite[Theorem~1.1]{WeakHessianCompanion}. Neither an
\(L^2\) condition on \(F\) nor an \(L^2\) Hessian of \(F\) is
required. The companion uses the same nonpositive Laplacian and
original test class as \eqref{eq:test-class}; its proof is independent
of the noncollapsed hypotheses and sectional-curvature assumption
used here. In the application below, the reference measure remains
\(\mathcal H^n\).

\subsection{Modified distance and the model equation}

Define the modified model distance by
\[
\mathfrak m_\kappa(r)=\int_0^r\sk(s)\dd s
=\begin{cases}
(1-\cos(\sqrt\kappa r))/\kappa,&\kappa>0,\\
r^2/2,&\kappa=0,\\
(\cosh(\sqrt{-\kappa}r)-1)/(-\kappa),&\kappa<0.
\end{cases}
\]
It obeys \(\mathfrak m_\kappa'=\sk\) and
\(\ck=1-\kappa\mathfrak m_\kappa\).

\begin{lemma}[A global bounded modification]
\label{weight:distance-modification}
Assume the conclusion of Theorem~\ref{frames:hessian}, and
write \(R_*=R_\kappa>0\) for its radius.
There are numbers \(0<a<b<R_*\),
depending only on \(\kappa\) and the chosen \(R_*\), such that for every \(p\in M\)
there are bounded globally Lipschitz \(F_p\) and bounded
continuous \(G_p\) satisfying \eqref{weight:hessian-assumption}
and
\begin{equation}
\label{weight:local-modification}
F_p=\mathfrak m_\kappa(d(p,\cdot)),\qquad
G_p=1-\kappa F_p\quad\text{on }B_a(p).
\end{equation}
For \(\kappa>0\), we may require
\(b<\pi/(2\sqrt\kappa)\).
\end{lemma}

\begin{proof}
Write \(r=d(p,\cdot)\) and \(q=r^2/2\), and fix \(0<a<b<R_*\).
Choose a smooth nonincreasing function \(\theta\colon[0,\infty)\to[0,1]\),
equal to one on \([0,a^2/2]\) and zero on
\([b^2/2,\infty)\), and set
\[
\chi(0)=0,\qquad
\chi'(q)=\theta(q)\frac{\sk(\sqrt{2q})}{\sqrt{2q}},\qquad
F_p=\chi(q).
\]
The quotient has a smooth extension at zero by its power
series. Consequently \(F_p\) is bounded and globally Lipschitz,
is constant outside \(B_b(p)\), and has the first value in
\eqref{weight:local-modification}. It need not lie in \(L^2(m)\);
boundedness and a bounded gradient are the hypotheses needed
for the weight.

For a compact test gradient \(Z=\nabla g\), the definition
\eqref{eq:weak-hessian} and the chain rule give
\begin{equation}
\label{weight:weak-chain}
H_{F_p}(Z,Z)(h)
=H_q(Z,Z)(h\chi'(q))
 +\int h\chi''(q)\langle\nabla q,Z\rangle^2\dd m.
\end{equation}
This identity involves only first derivatives of \(q\), so it
holds for the distributional Hessian. Local cutoffs keep its
first term inside \(B_{R_*}(p)\).
For \(r>0\), put
\(Z_r=\langle\nabla q,Z\rangle/r\); then
\(|Z_r|\le|Z|\), since the Lipschitz bound for \(r\) gives
\(|\nabla q|\le r\) almost everywhere.
The Hessian comparison of Theorem~\ref{frames:hessian} and
\(\chi'\ge0\) bound \eqref{weight:weak-chain} above by
\[
\int h\left[A(r)|Z_r|^2
                 +B(r)(|Z|^2-|Z_r|^2)\right]\dd m,
\]
where, after cancellation of the model-distance factors,
\begin{align*}
A(r)&=\chi'(r^2/2)+r^2\chi''(r^2/2)
     =\theta(r^2/2)\ck(r)+r\theta'(r^2/2)\sk(r),\\
B(r)&=\theta(r^2/2)\ck(r).
\end{align*}
The nonpositive singular part of \(H_q(Z,Z)\) remains
nonpositive on multiplication by \(\chi'\).
Both coefficients extend continuously with \(A(0)=B(0)=1\),
and vanish for \(r\ge b\). Since \(\theta'\le0\) and
\(\sk\ge0\) on \([0,b]\), we have \(A\le B\): flattening
the distance adds a nonpositive radial term. Hence the bounded
continuous function
\[
G_p(x)=\theta\bigl(r(x)^2/2\bigr)\ck\bigl(r(x)\bigr)
\]
gives \eqref{weight:hessian-assumption} globally.
On \(r<a\), it equals \(\ck(r)=1-\kappa F_p\), as required.
\end{proof}

\begin{theorem}[Alexandrov comparison]
\label{weight:alexandrov}
Under hypothesis \textup{(B)} of Theorem~\ref{main:equivalence},
\((M,d)\) is an \(n\)-dimensional Alexandrov space of curvature
at least \(\kappa\).
\end{theorem}

\begin{proof}
Theorems~\ref{heat:index} and \ref{frames:hessian} supply the
hypothesis of Lemma~\ref{weight:distance-modification}.
Apply Theorem~\ref{weight:every-geodesic} to \(F_p,G_p\).
For every geodesic \(\sigma\colon[0,1]\to B_a(p)\) of length
\(\ell\), \eqref{weight:local-modification} yields
\begin{equation}
\label{weight:model-inequality}
y''+\kappa\ell^2y\le\ell^2,
\qquad y(t)=\mathfrak m_\kappa(d(p,\sigma_t)).
\end{equation}

Consider a geodesic triangle \([pxy]\) of perimeter less than
\(a\), and let \(\sigma\) be its chosen side from \(x\) to
\(y\). The triangle inequality along the two routes through
the vertices gives \(d(p,\sigma_t)<a\), so
\eqref{weight:model-inequality} holds on this entire side.
The comparison triangle \([\bar p\bar x\bar y]\) in the
two-dimensional model of curvature \(\kappa\) exists at this
scale. For its corresponding side \(\bar\sigma\), the model
cosine law gives
\[
\bar y(t)=\mathfrak m_\kappa(d(\bar p,\bar\sigma_t)),\qquad
\bar y''+\kappa\ell^2\bar y=\ell^2,
\quad \bar y(0)=y(0),\quad\bar y(1)=y(1).
\]
This also holds for degenerate triangles by continuity.
Let \(h=y-\bar y\). It belongs to \(H^1_0(0,1)\), and
\(-h''-\kappa\ell^2h\ge0\) in distributions.
Testing with the negative part \(h_-\), by nonnegative
\(H^1_0\) approximation, gives
\[
\|(h_-)'\|_2^2\le\kappa\ell^2\|h_-\|_2^2.
\]
If \(\kappa\le0\), this forces \(h_-=0\). If \(\kappa>0\),
our scale choice gives \(\kappa\ell^2<\pi^2\), and the
Dirichlet Poincar\'e inequality gives the same conclusion.
Since \(\mathfrak m_\kappa\) is strictly increasing on the
distances under consideration,
\begin{equation}
\label{weight:point-side}
d(p,\sigma_t)\ge d(\bar p,\bar\sigma_t)
\qquad(0\le t\le1).
\end{equation}
Thus every sufficiently short triangle, with any chosen sides,
satisfies point-on-side comparison.

We spell out the localization for the final geometric step.
Fix \(o\in M\) and choose \(r>0\) with \(20r<a\).
The local converse to point-on-side comparison gives the
four-point \(\CBB(\kappa)\) inequality on \(B_r(o)\),
as follows. Point-on-side
comparison makes the comparison angle of a short hinge
nonincreasing in either radial parameter. Its limit defines
the angle, which is at least the comparison angle. The
adjacent-angle inequality follows from the model gluing lemma.
For a quadruple \(p,x,y,z\in B_r(o)\), take
\(w\in(p,z)\) with \(d(p,w)<r\). The radial legs from \(w\)
to these points have length at most \(3r\), so every triangle
used to compare points on these legs has perimeter at most
\(12r<a\). Write \(\widetilde\angle_\kappa(w;x,y)\) for
the comparison angle at \(w\) in the triangle \(wxy\).
For the actual angles at \(w\), insert the direction \(wp\) between
\(wx\) and \(wy\) in the angle triangle inequality. The resulting
terms pair with \(wz\) into two adjacent-angle sums, each at most
\(\pi\). Each comparison angle is at most its actual angle, so
\[
\widetilde\angle_\kappa(w;x,y)
+\widetilde\angle_\kappa(w;y,z)
+\widetilde\angle_\kappa(w;z,x)\le2\pi.
\]
Letting \(w\to p\), continuity of the comparison angles gives
the same inequality with \(w\) replaced by \(p\).
This is the local four-point inequality, obtained by localizing
the converse proof in
\cite[Theorem~8.14]{AKP2024}. All segments here are ambient
minimizing segments; no convexity of a metric ball is needed.
Hence \(M\) is locally \(\CBB(\kappa)\) by the local four-point
criterion in \cite[Theorem~8.30, condition~(1)]{AKP2024}.
Completeness and the length-space property allow the globalization theorem
\cite[Theorem~8.31]{AKP2024}, so \(M\) is \(\CBB(\kappa)\).

Finally, \(m=\mathcal H^n\) is finite on bounded sets and has
full support. Every nonempty bounded ball therefore has
positive finite \(n\)-dimensional Hausdorff measure, and hence
Hausdorff dimension \(n\). A countable exhaustion by such
balls shows \(\dim_H M=n\). This is the dimension of the
resulting Alexandrov space.
\end{proof}
